% संपूर्ण मराठी ग्रंथातील प्रकरण 33: द्वितीय-क्रम तर्कशास्त्राची अधिउपपत्ती
% हा संपादनयोग्य स्रोतखंड आहे. संकलनासाठी संपूर्ण स्रोतसंग्रहातील
% अक्षररूपे, पूर्वभाग, चिन्हव्याख्या आणि आकृती-स्रोत आवश्यक आहेत.
% मूळ: Open Logic Project; मजकूर आणि रूपांतर CC BY 4.0.
% यंत्रानुवाद, दुरुस्ती आणि तपासणी: OpenAI Codex —
% GPT-5.6 Sol आणि GPT-6 Sol, दोन्ही Ultra effort.
% दुरुस्ती-पुनरावलोकन: GPT-6.1 Sol, Ultra effort.
\chapter{द्वितीय-क्रम तर्कशास्त्राची अधिउपपत्ती}\label{sol:met:chap}









\section{परिचय}\label{sol:met:int:sec}


प्रथम-क्रम तर्कशास्त्राचे अनेक उपयुक्त गुणधर्म आहेत. ते निर्णेय नाही,
हे आपल्याला माहीत आहे; पण किमान ते स्वयंसिद्धकीकरणीय आहे. म्हणजे
प्रथम-क्रम तर्कशास्त्रासाठी निर्दोष आणि संपूर्ण अशा सिद्धता-प्रणाली
आहेत. त्या असा निष्पन्नता संबंध~$\Proves$ देतात की
वाक्यांचा कोणताही संच~$\Gamma$ आणि वाक्य~$\varphi$
यांसाठी $\Gamma \Entails \varphi$ जर आणि फक्त जर
$\Gamma \Proves \varphi$. विशेषतः, प्रथम-क्रम तर्कशास्त्रातील
वैध वाक्यांचा संच संगणनक्षमपणे प्रगणनीय आहे. या संगणनक्षम प्रगणनासाठी भाषा प्रभावी रीतीने
दिलेली असून सिद्धतांचेही प्रभावी प्रगणन करता येते, असे गृहीत आहे. एक संगणनक्षम फलन~$f\colon \Nat \to
\Sent[L]$ आहे, ज्यात~$f$ ची मूल्ये नेमकी~$\Lang{L}$ मधील सर्व
वैध वाक्ये आहेत. कारण निष्पत्त्यांचे प्रगणन
करता येते आणि त्यांपैकी एका~वाक्याला निष्पन्न
करणाऱ्या निष्पत्तींचे त्या वाक्य शी प्रतिचित्रण
करता येते. द्वितीय-क्रम तर्कशास्त्र प्रथम-क्रम तर्कशास्त्रापेक्षा
अधिक अभिव्यक्तिशील आहे; त्यामुळे त्याची वैधता ओळखणे
सर्वसाधारणपणे अधिक कठीण असते. प्रत्यक्षात, द्वितीय-क्रम
तर्कशास्त्र अनिर्णेय तर आहेच, पण त्याची वैधता
संगणनक्षमपणे प्रगणनीय सुद्धा नाही, हे आपण दाखवू. म्हणून
द्वितीय-क्रम तर्कशास्त्रासाठी निर्दोष, संपूर्ण आणि प्रभावी
सिद्धता-प्रणाली असू शकत नाही (मात्र निर्दोष पण अपूर्ण सिद्धता-प्रणाली उपलब्ध
आहेत आणि त्या संशोधनाचे महत्त्वाचे विषय आहेत).

प्रथम-क्रम तर्कशास्त्राचे आणखी दोन गुणधर्म आहेत: ते संहत
आहे (वाक्यांच्या संच~$\Gamma$ चा प्रत्येक सांत उपसंच
पूर्ततायोग्य असेल, तर $\Gamma$ स्वतःही पूर्ततायोग्य असतो)
आणि प्रगणनीय भाषेसाठी त्याला लोव्हेनहाइम--स्कोलेम प्रमेय लागू होते
($\Gamma$ चे अनंत प्रतिमान असेल, तर त्याचे
गणनीय अनंत प्रतिमानही असते). हे दोन्ही निष्कर्ष
द्वितीय-क्रम तर्कशास्त्रासाठी लागू होत नाहीत. कारण
द्वितीय-क्रम तर्कशास्त्राला प्रांताच्या आकाराविषयी
अशा गोष्टी व्यक्त करता येतात ज्या प्रथम-क्रम तर्कशास्त्राला
व्यक्त करता येत नाहीत.












\section{द्वितीय-क्रम अंकगणित}\label{sol:met:spa:sec}


पेआनो अंकगणिताची उपपत्ती~$\Th{PA}$ हिच्यात
$\Th{Q}$ ची आठ स्वयंसिद्धके समाविष्ट आहेत, हे आठवा:
\begin{align*}
& \lforall[x][\eq/[x'][\Obj 0]]\\
& \lforall[x][\lforall[y][(\eq[x'][y'] \lif \eq[x][y])]]\\
& \lforall[x][(\eq[x][\Obj 0] \lor \lexists[y][\eq[x][y']])]\\
& \lforall[x][\eq[(x + \Obj 0)][x]]\\
& \lforall[x][\lforall[y][\eq[(x + y')][(x + y)']]]\\
& \lforall[x][\eq[(x \times \Obj 0)][\Obj 0]]\\
& \lforall[x][\lforall[y][\eq[(x \times y')][((x \times y) + x)]]]\\
& \lforall[x][\lforall[y][(x < y \liff \lexists[z][\eq[(z' + x)][y]])]]\\
\intertext{आणि पुढील स्वरूपाची सर्व वाक्येही:}
& (\varphi(\Obj 0) \land \lforall[x][(\varphi(x) \lif \varphi(x'))]) \lif \lforall[x][\varphi(x)].
\end{align*}
शेवटचे हे एक ``रूपबंध'' आहे: अंकगणिताच्या भाषेतील
अमर्याद वाक्ये निर्माण करणारा नमुना, प्रत्येक
सूत्र~$\varphi(x)$ साठी एक वाक्य. इतर मुक्त चल असतील,
तर आरंभी त्यांच्यावर सार्विक संख्यापक जोडून वाक्यरूप घेतो. याला (प्रथम-क्रम)
\emph{विगमनाचा स्वयंसिद्धक-रूपबंध} म्हणतो.
\emph{द्वितीय-क्रम} पेआनो अंकगणित~$\Th{PA^2}$ मध्ये
विगमन एकाच वाक्यात मांडता येते. $\Th{PA^2}$ मध्ये
वरील पहिली आठ स्वयंसिद्धके आणि (द्वितीय-क्रम)
\emph{विगमनाचे स्वयंसिद्धक} असते:
\[\lforall[X][((X(\Obj 0) \land \lforall[x][(X(x) \lif X(x'))]) \lif \lforall[x][X(x)])].
\]
याचा अर्थ असा: प्रांताचा उपसंच~$X$ यामध्ये
$\Assign{\Obj{0}}{M}$ असेल आणि कोणत्याही~$x \in \Domain{M}$
बरोबर~$\Assign{\prime}{M}(x)$ ही असेल (म्हणजे तो
``उत्तराधिकारी क्रियेखाली बंद'' असेल), तर प्रांत मधील
सर्व घटक त्यात असतील (म्हणजे $X =
\Domain{M})$.

विगमनाचे स्वयंसिद्धक असा निर्वाळा देते की ते पूर्ण करणाऱ्या
कोणत्याही संरचनेच्या~$\Domain{M}$ मध्ये
स्वयंसिद्धकांना आवश्यक असलेलेच घटक असतात;
म्हणजे $n \in \Nat$ साठी~$\num{n}$ ची मूल्ये.
$\Th{PA^2}$ च्या प्रतिमानात अमानक संख्या नसतात.

\begin{thm}
\label{sol:met:spa:thm:sol-pa-standard}
जर $\Sat{M}{\Th{PA^2}}$ असेल, तर $\Domain{M} =
\Setabs{\Value{\num{n}}{M}}{n \in \Nat}$.
\end{thm}

\begin{proof}
$N = \Setabs{\Value{\num{n}}{M}}{n \in \Nat}$ असे मानू
आणि $\Sat{M}{\Th{PA^2}}$ असे गृहीत धरू. अर्थात,
कोणत्याही $n \in \Nat$ साठी
$\Value{\num{n}}{M} \in \Domain{M}$; म्हणून
$N \subseteq \Domain{M}$.

आता उलट दिशेचा अंतर्भाव दाखवू. $s(X) = N$ असे
चल-मूल्यांकन $s$ घेऊ. गृहीतकानुसार,
\begin{align*}
\Sat{M}{\lforall[X][((X(\Obj 0) \land \lforall[x][(X(x)
      \lif X(x'))]) \lif \lforall[x][X(x)])]}, & \text{ म्हणून}\\
\Sat{M}{(X(\Obj 0) \land \lforall[x][(X(x) \lif X(x'))]) \lif
  \lforall[x][X(x)]}[s]. &
\end{align*}
या सशर्त विधानाचे पूर्वांग पाहू. $\Value{\Obj{0}}{M} \in
N$, म्हणून $\Sat{M}{X(\Obj{0})}[s]$. दुसरे संयुक्तपद,
$\lforall[x][(X(x) \lif X(x'))]$, हेही पूर्ण होते.
कारण $x \in N$ असे समजू. $N$ च्या व्याख्येनुसार,
काही~$n$ साठी $x = \Value{\num{n}}{M}$. त्यावरून
$\Assign{\prime}{M}(x) = \Value{\num{n+1}}{M} \in N$.
म्हणून $\Assign{\prime}{M}(x) \in N$.

त्यामुळे $\Sat{M}{X(\Obj 0) \land \lforall[x][(X(x) \lif X(x'))]}[s]$.
म्हणून $\Sat{M}{\lforall[x][X(x)]}[s]$. याचा अर्थ,
प्रत्येक $x \in \Domain{M}$ साठी $x \in s(X) = N$.
म्हणून $\Domain{M} \subseteq N$.
\end{proof}

\begin{cor}
\label{sol:met:spa:cor:sol-pa-categorical}
$\Th{PA^2}$ ची कोणतीही दोन प्रतिमाने समरूप असतात.
\end{cor}

\begin{proof}
\ref{sol:met:spa:thm:sol-pa-standard} नुसार,
$\Th{PA^2}$ च्या कोणत्याही प्रतिमानाच्या प्रांतात
$\Value{\num{n}}{M}$ ही मूल्येच असतात. असे प्रत्येक
प्रतिमान~$\Th{Q}$ चेही प्रतिमान असते.
\ref{mod:mar:stm:prop:thq-standard} नुसार
ते मानक असते, म्हणजे~$\Struct{N}$ शी समरूप असते.
\end{proof}

वर $\Th{PA^2}$ ची व्याख्या पहिली आठ अंकगणितीय स्वयंसिद्धके
आणि द्वितीय-क्रम विगमनाचे स्वयंसिद्धक असलेली उपपत्ती म्हणून
केली. प्रत्यक्षात, द्वितीय-क्रम तर्कशास्त्राच्या
अभिव्यक्तीक्षमतेमुळे द्वितीय-क्रम पेआनो अंकगणितासाठी
अंकगणितीय स्वयंसिद्धकांपैकी फक्त \emph{पहिली दोन}
आणि विगमन पुरेसे असते.

\begin{prop}\label{sol:met:spa:prop:sol-pa-definable}
पहिली दोन अंकगणितीय स्वयंसिद्धके (उत्तराधिकारी स्वयंसिद्धके)
आणि द्वितीय-क्रम विगमनाचे स्वयंसिद्धक असलेली द्वितीय-क्रम
उपपत्ती $\Th{PA^{2\dagger}}$ असो. मग $\le$, $+$ आणि
$\times$ हे $\Th{PA^{2\dagger}}$ मध्ये परिभाष्य आहेत.
\end{prop}

\begin{proof}
$\le$ परिभाष्य आहे हे दाखवण्यासाठी असे
सूत्र~$\varphi_\le(x, y)$ शोधावे लागेल की
$\Sat{N}{\varphi_\le(\num n, \num m)}$ जर आणि फक्त जर
$n \le m$. पुढील सूत्र पाहू:
\[\psi(x, Y) \ident Y(x) \land \lforall[y][(Y(y) \lif Y(y'))]
\]
संच $Y \subseteq \Nat$ कडून $\psi(\num n, Y)$ ची
पूर्तता झाल्यास $\Setabs{m}{n \le m} \subseteq Y$ असते.
उलट, फक्त हा अंतर्भाव पुरेसा नाही; Y उत्तराधिकारी क्रियेखाली
बंदही असावा लागतो. n पासून पुढच्या सर्व नैसर्गिक संख्यांचा
संच स्वतः वरील सूत्राची पूर्तता करतो आणि अशा प्रत्येक Y मध्ये
तो समाविष्ट असतो. म्हणून
$\varphi_\le(x, y) \ident
\lforall[Y][(\psi(x, Y) \lif Y(y))]$ असे घेता येते.

बेरीज परिभाष्य आहे हे पाहण्यासाठी लक्षात घ्या की $k+l=m$
जर आणि फक्त जर असे फलन $u$ असेल की $u(0)=k$,
प्रत्येक $n$ साठी $u(n')=u(n)'$ आणि $m = u(l)$.
या समतुल्यतेने $\Th{PA^{2\dagger}}$ मध्ये पुढील
सूत्राद्वारे बेरीज परिभाषित करता येते:
\[\varphi_+(x,y,z) \ident \lexists[u][(u(\Obj 0)=x \land \lforall[w][u(w')=u
 (w)'] \land u(y)=z)]
\]
स्पष्टच, $\Sat{N}{\varphi_+(\num k, \num l, \num m)}$
जर आणि फक्त जर $k+l=m$.
\end{proof}

\begin{prob}
\ref{sol:met:spa:prop:sol-pa-definable} चा पुरावा पूर्ण करा.
\end{prob}












\section{द्वितीय-क्रम तर्कशास्त्र स्वयंसिद्धकीकरणीय नाही}\label{sol:met:nax:sec}


\begin{thm}
\label{sol:met:nax:thm:sol-undecidable}
द्वितीय-क्रम तर्कशास्त्र अनिर्णेय आहे.
\end{thm}

\begin{proof}
प्रथम-क्रम वाक्य प्रथम-क्रम तर्कशास्त्रात वैध असेल,
जर आणि फक्त जर ते द्वितीय-क्रम तर्कशास्त्रात वैध असेल.
प्रथम-क्रम तर्कशास्त्र अनिर्णेय आहे.
\end{proof}

\begin{thm}
\label{sol:met:nax:cor:sol-not-axiomatizable}
द्वितीय-क्रम तर्कशास्त्राच्या मानक चिन्हार्थमीमांसेसाठी
निर्दोष, संपूर्ण आणि प्रभावी निष्पत्ती प्रणाली अस्तित्वात नाही.
\end{thm}

\begin{proof}

अंकगणिताच्या प्रथम-क्रम भाषेतील वाक्य~$\varphi$ असो.
$\Sat{N}{\varphi}$ जर आणि फक्त जर $\Th{PA^2} \Entails \varphi$.
$\mathsf{P}$ हा~$\Th{PA^2}$ च्या नऊ स्वयंसिद्धकांचा संयोग असो.
$\Th{PA^2} \Entails \varphi$ जर आणि फक्त जर
$\Entails \mathsf{P} \lif \varphi$; म्हणजे प्रत्येक रचनेसाठी
$\Sat{M}{\mathsf{P} \lif \varphi}$. आता पुढील वाक्य पाहू:
$\lforall[z][\lforall[u][\lforall[u'][\lforall[u''][\lforall[L][(\mathsf{P}' \lif \varphi')]]]]]$.
हे $\Obj{0}$ च्या जागी $z$, $\prime$ च्या जागी
एकस्थानी फलनचल $u$, $+$ आणि $\times$ यांच्या जागी
अनुक्रमे द्विस्थानी फलनचले $u'$ आणि $u''$,
आणि $<$ च्या जागी द्विस्थानी संबंधचल~$L$ ठेवून
नंतर सार्वत्रिक संख्यापनाने मिळते. येथे नव्याने वापरलेले
सर्व चले आधीच्या सूत्रांत नसलेले निवडतो; आवश्यक असल्यास
आधी बद्ध चलांची नावे बदलतो, म्हणून capture होत नाही. हे शुद्ध द्वितीय-क्रम
तर्कशास्त्रातील वैध वाक्य असेल, जर आणि फक्त जर मूळ
वाक्य वैध असेल, जर आणि फक्त जर
$\Th{PA^2} \Entails \varphi$, जर आणि फक्त जर
$\Sat{N}{\varphi}$. म्हणून द्वितीय-क्रम तर्कशास्त्रासाठी
निर्दोष, संपूर्ण आणि प्रभावी सिद्धता-प्रणाली असती, तर
प्रथम-क्रम अंकगणितीय वाक्यांचे आणि त्यांच्या वरील निर्मित
वाक्यांसाठीच्या सिद्धतांचे समांतर प्रभावी प्रगणन वापरून
$\Struct{N}$ मध्ये सत्य असलेल्या वाक्यांचे
संगणनक्षम प्रगणन $f\colon \Nat \to \Sent[L_A]$
परिभाषित करता आले असते. हे फलन काही प्रथम-क्रम
सूत्र~$\psi_f(x, y)$ द्वारे~$\Th{Q}$ मध्ये निरूपित
करता आले असते. मग सूत्र~$\lexists[x][\psi_f(x, y)]$
हे~$\Struct{N}$ मधील सत्य प्रथम-क्रम वाक्यांचा संच
परिभाषित करील. हा टार्स्कीच्या प्रमेयाशी विरोध आहे.
\end{proof}














\section{द्वितीय-क्रम तर्कशास्त्र संहत नाही}\label{sol:met:com:sec}

\begin{explain}
वाक्यांचा संच~$\Gamma$ याचा प्रत्येक सांत उपसंच
पूर्ततायोग्य असेल, तर त्याला \emph{सांततः पूर्ततायोग्य}
म्हणतो. प्रथम-क्रम तर्कशास्त्राचा असा गुणधर्म आहे की
वाक्यांचा संच~$\Gamma$ सांततः पूर्ततायोग्य
असेल, तर तो पूर्ततायोग्य असतो. या गुणधर्माला
\emph{संहतता} म्हणतात. याचे परिणाम-संबंध वापरणारे
समतुल्य रूप असे: $\Gamma \Entails \varphi$ असेल,
तर काही सांत उपसंच~$\Gamma_0 \subseteq \Gamma$
यासाठी आधीच $\Gamma_0 \Entails
\varphi$ असते. या रूपात हा संपूर्णता प्रमेयाचा थेट
परिणाम आहे: $\Gamma \Entails \varphi$ असेल, तर
संपूर्णतेमुळे $\Gamma \Proves \varphi$. पण
निष्पत्तीमध्ये~$\Gamma$ मधील सांत
संख्येनेच वाक्ये वापरता येतात.

द्वितीय-क्रम तर्कशास्त्रासाठी संहतता लागू होत नाही.
द्वितीय-क्रम वाक्यांचे असे संच आहेत की ते
सांततः पूर्ततायोग्य असतात, पण पूर्ततायोग्य नसतात;
आणि ते एखाद्या~$\varphi$ चे परिणाम देतात, पण त्यांचा
कोणताही सांत उपसंच त्या~$\varphi$ चे परिणाम देत नाही.
\end{explain}


\begin{thm}
\label{sol:met:com:thm:sol-not-compact}
द्वितीय-क्रम तर्कशास्त्र संहत नाही.
\end{thm}

\begin{proof}
हे आठवा:
\[\fn{Inf} \ident \lexists[u][(\lforall[x][\lforall[y][(\eq[u(x)][u(y)]
      \lif \eq[x][y])]] \land \lexists[y][\lforall[x][\eq/[y][u(x)]]])]
\]
संरचनेचा प्रांत अनंत असेल, जर आणि फक्त जर
तिच्यात या सूत्राची पूर्ति होते. प्रांतात किमान
$n$ घटक आहेत, असे सांगणारे वाक्य
$\varphi^{\ge n}$ असो; उदाहरणार्थ,
\[\varphi^{\ge n} \ident \lexists[x_1][\dots\lexists[x_n][(\eq/[x_1][x_2]
    \land \eq/[x_1][x_3] \land \dots \land \eq/[x_{n-1}][x_n])]].
\]
वाक्यांचा पुढील संच पाहू:
\[\Gamma = \{\lnot \fn{Inf}, \varphi^{\ge 1}, \varphi^{\ge 2}, \varphi^{\ge 3}, \dots\}.
\]
तो सांततः पूर्ततायोग्य आहे: कोणताही सांत उपसंच~$\Gamma_0
\subseteq \Gamma$ घेतल्यास असा $k$ निवडता येतो की
$\varphi^{\ge k} \in \Gamma$ असते, पण $n>k$ साठी
कोणतेही $\varphi^{\ge n} \in \Gamma_0$ नसते.
$\Domain{M}$ मध्ये $k$ घटक असतील, तर
$\Sat{M}{\Gamma_0}$. मात्र $\Gamma$ पूर्ततायोग्य
नाही: $\Sat{M}{\lnot \fn{Inf}}$ असेल, तर
$\Domain{M}$ सांत असला पाहिजे; समजा त्याचा
आकार~$k$ आहे. मग $\Sat/{M}{\varphi^{\ge k+1}}$.
\end{proof}

\begin{prob}
संच~$\Gamma$ आणि वाक्य~$\varphi$ यांचे असे
उदाहरण द्या की $\Gamma \Entails \varphi$, पण प्रत्येक
सांत उपसंच~$\Gamma_0 \subseteq
\Gamma$ साठी $\Gamma_0 \Entails/ \varphi$.
\end{prob}













\section{द्वितीय-क्रम तर्कशास्त्रात लोव्हेनहाइम--स्कोलेम प्रमेय लागू होत नाही}\label{sol:met:lst:sec}

\begin{explain}
(अधोगामी) लोव्हेनहाइम--स्कोलेम प्रमेयानुसार, अनंत
प्रतिमान असलेल्या वाक्यांच्या प्रत्येक संचाला
गणनीय प्रतिमानही असते. हे प्रमेयदेखील
संपूर्णता प्रमेयाचा परिणाम आहे: संपूर्णतेचा पुरावा
वाक्यांच्या कोणत्याही सुसंगत संचासाठी प्रतिमान
तयार करतो आणि ते प्रतिमान गणनीय असते.
ऊर्ध्वगामी लोव्हेनहाइम--स्कोलेम प्रमेयही आहे.
त्यानुसार, वाक्यांच्या संचाला
गणनीय अनंत प्रतिमान असेल, तर त्याला
अगणनीय प्रतिमानही असते. ही दोन्ही
प्रमेये द्वितीय-क्रम तर्कशास्त्रात लागू होत नाहीत.
\end{explain}


\begin{thm}
\label{sol:met:lst:thm:sol-no-ls} लोव्हेनहाइम--स्कोलेम प्रमेय
द्वितीय-क्रम तर्कशास्त्रात लागू होत नाही:
काही वाक्यांना अनंत प्रतिमाने असतात,
पण गणनीय प्रतिमाने नसतात.
\end{thm}

\begin{proof}
हे आठवा:
\[\fn{Count} \ident \lexists[z][\lexists[u][\lforall[X][((X(z) \land
      \lforall[x][(X(x) \lif X(u(x)))]) \lif \lforall[x][X(x)])]]]
\]
हे संरचना~$\Struct{M}$ मध्ये सत्य असेल,
जर आणि फक्त जर $\Domain{M}$ गणनीय असेल.
म्हणून $\lnot\fn{Count}$ हे~$\Struct{M}$ मध्ये
सत्य असेल, जर आणि फक्त जर $\Domain{M}$
अगणनीय असेल. अशा संरचना
आहेत---कोणताही अगणनीय संच प्रांत
म्हणून घ्या; उदाहरणार्थ, $\Pow{\Nat}$ किंवा~$\Real$.
म्हणून $\lnot \fn{Count}$ ला अनंत प्रतिमाने
आहेत, पण गणनीय प्रतिमाने नाहीत.
\end{proof}

\begin{thm}
काही वाक्यांना गणनीय अनंत प्रतिमाने
असतात, पण अगणनीय प्रतिमाने नसतात.
\end{thm}

\begin{proof}
$\fn{Count} \land \fn{Inf}$ हे $\Nat$ मध्ये
सत्य असते; पण $\Domain{M}$ अगणनीय
असलेल्या कोणत्याही संरचना~$\Struct{M}$
मध्ये ते सत्य नसते.
\end{proof}




